\documentclass[journal]{IEEEtran}

\usepackage{cite}
\usepackage{amsmath}
\usepackage{graphicx}
\usepackage{url}
\usepackage{hyperref}
\usepackage{booktabs}
\usepackage{listings}

\lstset{
  basicstyle=\ttfamily\small,
  breaklines=true,
  frame=single,
  xleftmargin=1em,
  framexleftmargin=0.5em
}

\begin{document}

\title{Reanimating the First AI: The Logic Theorist Runs Again in IPL-V}

\author{Jeff~Shrager,~Ph.D.\\
Bennu Climate, Inc.\\
\texttt{jeff@shrager.org}}

\maketitle

\begin{abstract} >>> this is very good -- note minor detail re Lisp <<
The Logic Theorist (LT), created by Allen Newell, J.C.\ (Cliff) Shaw,
and Herbert Simon at the RAND Corporation and Carnegie Tech in
1955--1956, is widely regarded as the first artificial intelligence
program. LT was implemented in the Information Processing Language
(IPL), a list-processing language that directly foreshadowed Lisp.
>>> IPL was an explicit precursor of lisp per McCarthy <<< Unlike
Lisp, however, IPL has been essentially extinct for over half a
century; no implementations are known to have survived, and no
community of practitioners remains. Here I describe the construction
of a new IPL-V interpreter, written in Common Lisp, and the faithful
reanimation of the original Logic Theorist from code transcribed
directly from Stefferud's 1963 RAND technical report. The reanimated
LT successfully proves 16 of 23 theorems from Chapter~2 of Whitehead
and Russell's \textit{Principia Mathematica}, results that are
historically consistent with the original. This effort continues a
program of reanimating foundational AI systems on their original
platforms, following the recent restoration of Weizenbaum's ELIZA on a
reconstructed CTSS. A distinctive aspect of this project was the
essential role played by large language models (LLMs) as collaborative
debugging partners---learning an essentially dead programming language
alongside the author, and ultimately proving indispensable in
navigating thousands of lines of execution traces from a heuristic
program whose correctness could not easily be verified.
\end{abstract}

\begin{IEEEkeywords}
Logic Theorist, IPL-V, artificial intelligence history, software
archaeology, Newell, Simon, Shaw, list processing, digital preservation
\end{IEEEkeywords}

%% ============================================================
\section{Introduction}
\label{sec:intro}
%% ============================================================

In the summer of 1956, a small group of researchers gathered at
Dartmouth College for a workshop on what John McCarthy had dubbed
``artificial intelligence.'' Among them were Allen Newell and Herbert
Simon, who arrived with something none of the other attendees had: a
working program that could discover proofs of mathematical theorems.

>>> A clarification: The above makes it seem like what was special
about ehat Simon and Newell brought was that is was a theorem prover.
But there were no other working AIs *at all* at the workshop (as far
as I know), so what Simon and Newell brought was a working AI -- that
was what was impressive -- it happened to be a theorem prover. <<<

The Logic Theorist (LT) had already proved 38 of the first 52 theorems
in Chapter~2 of Whitehead and Russell's \textit{Principia
  Mathematica}~\cite{russell1910principia}, in some cases finding
proofs more elegant than those produced by the human authors. It was,
arguably, the first artificial intelligence. >>>Take the Samuels'
content below and make it a (shorter) footnote here instead.<<<

LT was implemented in the Information Processing Language (IPL), a
language created by the same team---Newell, Shaw, and
Simon---specifically for building cognitive
models~\cite{newell1961ipl}. IPL introduced list manipulation, dynamic
memory allocation, recursion, and >>> Add: symbols as first-class
objects <<< higher-order functions. In a meaningful sense, >>> ``in a
meaningful sense'' is not a useful phrase. It was EXPLICITLY (per
McCarthy!) a direct ancenstor of Lisp!) <<< it was the direct ancestor
of Lisp~\cite{mccarthy1960lisp}. Yet while Lisp has thrived for nearly
seven decades, IPL has been extinct since the mid-1960s. No
implementations are known to survive, and to the author's knowledge,
no one has previously attempted to build an IPL-V interpreter and run
the original AI programs.

>>>>>>>>>>>>>>>>>>>>>>>>>>>>>>>>>>>>>>>>>>>>>>>>>>>>>>>>>>>

Re above: Many IPL-V implementations had been around and proposed in
the 1960s, but not since. I know of one by Howard Richman
\cite{richman1995epam} Richman wrote to me (personal communication,
Sep 26 2025, quoted with permission of HR): ``I just interpreted IPL-V
in Microsoft Basic using the IPL-V manual that I found at the CMU
library.''

<<<<<<<<<<<<<<<<<<<<<<<<<<<<<<<<<<<<<<<<<<<<<<<<<<<<<<<<<<<

This paper describes the construction of a complete IPL-V interpreter
in Common Lisp, and the faithful reanimation of the Logic Theorist
from code transcribed directly from Stefferud's 1963 RAND technical
report~\cite{stefferud1963lt}. More than 99\% of the LT code running
in this reconstruction is the original code from that report. The
reanimated LT proves 16 of 23 attempted theorems from
\textit{Principia Mathematica}, with results that are historically
consistent with the original system's known behavior.

This work continues a program of reanimating foundational AI and
cognitive science programs on their original platforms. In prior work,
a team including the present author restored Weizenbaum's original
ELIZA to operation on a reconstructed CTSS running on an emulated IBM
7094~\cite{lane2025eliza}. The present project differs in both method
and character: where the ELIZA reanimation reconstructed an entire
hardware and operating system stack, the LT reanimation emulates the
abstract IPL-V machine in its direct descendant language, Common Lisp.
And where ELIZA was a team effort, the LT project was largely a solo
endeavor---albeit one in which large language models played an
unexpectedly essential role as collaborative debugging partners.

%% ============================================================
\section{Historical Background}
\label{sec:history}
%% ============================================================

\subsection{The Origins of the Logic Theorist}

The story of the Logic Theorist begins in the early 1950s at the RAND
Corporation in Santa Monica, California. Herbert Simon, a political
scientist consulting at RAND, saw a printer rendering a map using
ordinary characters and realized that a machine capable of
manipulating symbols could simulate decision-making---and perhaps
thought itself. Allen Newell, a RAND scientist studying logistics and
organization theory, had a complementary epiphany watching Oliver
Selfridge present his work on pattern recognition in 1954: he saw how
the interaction of simple programmable units could produce complex,
intelligent behavior.

Newell and Simon began collaborating on a program that could prove
theorems in propositional logic, enlisting RAND programmer J.C.\
(Cliff) Shaw to handle the implementation. The first version was
famously hand-simulated: Simon assembled his wife and children, gave
each a card representing a component of the program, and had them
execute the algorithm by passing cards back and forth.

By 1956, the Logic Theorist was running on the
JOHNNIAC~\cite{gruenberger1968johnniac}, one of the first
Turing--Von~Neumann machines, built at RAND and named after Von
Neumann himself. The JOHNNIAC had just 4,096 40-bit words of
memory. Within these severe constraints, LT proved 38 of the first 52
theorems in Chapter~2 of \textit{Principia Mathematica}, discovering
proofs by heuristic search through a space of possible proof steps.
For one theorem, LT found a proof more elegant than the one published
by Whitehead and Russell.

It is worth noting that Art Samuel's checkers-playing program existed
a couple of years before LT, and it too was heuristic. But Samuel's
goal was to build a checkers player, whereas Newell, Simon, and Shaw's
explicit goal was to model human cognition. They held specific
hypotheses about how humans reason and about the ``physical symbol
system'' that, by their hypothesis, underlies intelligence. Their
research program spanned many domains over the following decades. This
is not to diminish Samuel's brilliant effort---he also built arguably
the first machine learning program, and the first system to learn from
self-play.

\subsection{IPL: The First List-Processing Language}

To implement LT, Newell, Shaw, and Simon needed a programming language
suited to symbolic reasoning. No such language existed, so they
created one: the Information Processing Language~\cite{newell1961ipl}.
IPL went through several versions; IPL-V, released in 1964, was the
only version made publicly available~\cite{newell1964iplv}.

IPL-V runs on an abstract register machine with a push-down stack,
working registers (W0--W9), >>> this isn't quite right. These aren't
really ``registers'', but cells, and any cell can be pushed and
popped, so it's not like there was ``a (single) push down stack'' <<<
communication registers (H0--H5) >>> These are control cells, not
communication cells (some are explicitly for communication, but some
are not -- like H3 which is the cycle counter, so ``control'' is a
more general term<<<, and a symbol table of cells >> This minimizes
the concept of symbol in IPL-V -- symbols were central, and every cell
had a symbol associated with it!) <<<. Each cell has three fields: a
symbol (SYMB), a link (LINK), and a set of process qualifiers
(PQ). >>> The Symb/Link/PQ is too much detail <<< Everything in
IPL-V---data, routines, even the program itself---is a list of
symbols. Routines are invoked by executing their name as an
instruction; the machine traverses the list associated with that name,
interpreting each cell as an instruction.

The language's built-in operations are called J-functions (because
their names begin with ``J'': J0, J1, J2, and so on). There are
approximately 150 of these, implementing operations from basic list
manipulation to generator protocols (a form of lazy iteration that
anticipates modern iterators). In the original IPL-V, many J-functions
were themselves written in IPL, although some required machine-level
primitives.

IPL is, in a meaningful sense, >>> You need to stop using this phrase
-- and see this is redundant from above <<< the direct ancestor of
Lisp~\cite{mccarthy1960lisp}. It introduced list and symbol
manipulation, the idea that programs are themselves lists, dynamic
memory allocation, recursion, and higher-order functions. Having now
spent months implementing an IPL-V interpreter, I can report that IPL
had essentially \textit{everything} that Lisp has, with one critical
exception: homoiconic syntax. >>> Not exactly, esp. if you count the
other important thing of GC! <<< Where Lisp unified code and data in a
single, elegant, parenthesized notation, IPL expressed its ideas in an
assembly-language style that, while powerful, was far more difficult
to read and write. McCarthy's genius was not in inventing list
processing \textit{de novo}, but in finding the syntactic form that
made these ideas accessible. Having programmed in Lisp for fifty years
and attributed its brilliance entirely to McCarthy, building this
interpreter gave me a deepened appreciation for Newell, Simon, and
Shaw's foundational contributions.

\subsection{The McCorduck Interviews} >>> Should this really be a heading -- instead these should be folded in at the right places <<<

Among the historical materials that informed this project were
transcripts of interviews conducted by Pamela McCorduck for her
pioneering history of AI~\cite{mccorduck1979machines}. These
interviews, originally obtained for the ELIZA reanimation
project~\cite{lane2025eliza}, include conversations with Newell,
Simon, and Shaw themselves. One passage proved particularly resonant:
Shaw describes the difficulty of debugging a heuristic program---one
that does not crash, does not produce obviously wrong output, but
simply generates vast quantities of behavior that might or might not
be correct. As described in Section~\ref{sec:debugging}, I encountered
precisely this problem, sixty years later, with the same program.

>>>>>>>>>>>>>>>>>>>>>>>>>>>>>>>>>>>>>>>>>>>>>>>>>>>>>>>>>>>
Here's the quote from Shaw's interview with McCorduck:

``When we first used dynamic storage allocation and it kept the list
of available space and dumped unused space back on this list, reused
space back on this list so it could be available, the memory was
churning over in a way that none of us had ever experienced
before. You take a memory dump and it was completely laced through
with these links from these lists. So naturally we didn't use memory
dumps very much except when we were forced to by the machine
failure. Rather we would use debugging tools that would retrieve lists
and dump out their contents both in the process of the program and
also after something had failed. But initially when we started using
list processing, available space lists, the symptom of trouble was
inevitably out of available space. All kinds of errors ended up being
out of available space because they would destroy the available space
list of in some way use it up. Well finally we got that fixed to where
that was not the immediate response with programming errors and later
on, particularly in chess, it got so that the machine would produce
behavior but you weren't at all sure that what you had made in the way
of changes for this last run actually influenced that behavior in the
way they were intended to. It would behave reasonably in some
sense. It would go through the motions of making its move and wait for
yours but it was difficult to debug in a completely different
sense. Not that it fell apart but that it wouldn't fall apart. You
weren't sure that new programs, revisions, had the effect that you
intended because it would do something reasonable -- it wouldn't
always do what you predicted and then you were in a quandry as to
whether you had a bug or whether this was a consequence of what you
had specified.'' (conversation with J. C. Shaw on June 16, 1975)

<<<<<<<<<<<<<<<<<<<<<<<<<<<<<<<<<<<<<<<<<<<<<<<<<<<<<<<<<<<


%% ============================================================
\section{Source Materials and Transcription}
\label{sec:sources}
%% ============================================================

Unlike the ELIZA reanimation, which began with the dramatic discovery
of original source code printouts in Weizenbaum's archives at
MIT~\cite{lane2025eliza}, the source materials for the Logic Theorist
were already publicly available, if not widely known. Two documents
were indispensable:

\begin{enumerate}
\item \textbf{Stefferud (1963)}, \textit{The Logic Theory Machine: A
  Model Heuristic Program}, RAND Corporation memorandum
  RM-3731~\cite{stefferud1963lt}. This report contains a complete
  listing of the LT program in IPL-V, along with detailed descriptions
  of its M-routines (the main heuristic procedures).

\item \textbf{Newell et al.\ (1964)}, \textit{Information Processing
  Language-V Manual}, 2nd edition~\cite{newell1964iplv}. The
  definitive reference for IPL-V, including the semantics of all
  J-functions, cell structure conventions, and the generator protocol.
\end{enumerate}

The IPL-V manual proved to be an exceptionally well-written
document. >>> Both Stefferud's paper and the IPL-V manual were
exceptionally well written! <<< There were, however, occasional
ambiguities---places where the specification assumed familiarity with
conventions of 1960s computing (such as BCD character encoding) that
are no longer common knowledge.

The LT code was transcribed from Stefferud's report into a Google
spreadsheet by the author and Anthony Hay. Subsequently, Rupert Lane
identified several transcription errors using his ``gridlock''
tool~\cite{lane2024gridlock}, software designed to accurately extract
tabular code from PDF documents. Despite multiple people putting
multiple sets of eyes on the code, a surprising number of small
transcription errors persisted through several rounds of review.

The transcribed code was converted into an S-expression format
(\texttt{.liplv} files) suitable for loading by the Lisp-based
interpreter. Conveniently, Lisp comment characters work in these
files, allowing corrections and annotations to be documented inline.

%% ============================================================
\section{The IPL-V Interpreter}
\label{sec:interpreter}
%% ============================================================

Rather than reconstructing >>> emulating, not reconstructing <<< the
original hardware environment (as was done for ELIZA on the IBM 7094
and CTSS) >>>citation to Lane et al<<, I chose to emulate the IPL-V
abstract machine in Common Lisp. This decision was driven by several
considerations. First, the JOHNNIAC on which LT originally ran no
longer exists and, to my knowledge, no emulator for it is
available. Second, IPL-V is defined as an abstract machine
specification, not as a program for a specific computer; it was
implemented on many different machines through the mid-1960s before
being supplanted by Lisp. Third, there is a pleasing circularity in
implementing IPL-V---Lisp's direct ancestor---in Lisp itself.

>>> There an IPL-V interpreter that comes with B7094 that It comes
with a number of emulator demo scripts including two for IPLV The
interpreter is a binary program running under a variant of the IBSYS
batch operating system. So no source code for IPL nor any
documentation on how to use it (eg does it have any run time
switches): https://www.cozx.com/dpitts/ibm7090.html I tried running LT
on this, but it failed very early, and beause there is no source, it
is essentially un-debuggable. Moreover -- and here's the important
points -- I wanted to write my own IPL because I wanted a deep
understanding of IPL in addition to that of LT and GPS and other code
running on it. The goal wasn't just to get it to run, but to elucidate
it, and the infrastructure of the time. <<<

My interpreter (\texttt{iplv.lisp}) implements the complete IPL-V
abstract machine as described in the 1964 manual: the cell table with
SYMB, LINK, and PQ fields; the communication registers H0--H5 (where
H0 serves as the parameter/result register, H1 as the program counter
with its associated call stack, and H2 as the free list); the working
registers W0--W9; and the push-down stack. Lists are stored as linked
chains of cells, exactly as specified.

The core of the interpreter is a function \texttt{ipl-eval}, which
fetches and executes instructions by traversing the list named in the
program counter (H1). \texttt{ipl-eval} is re-entrant, supporting the
recursive subroutine calls that IPL-V requires. The J-functions are
implemented as Lisp functions dispatched by name. I implemented only
those J-functions actually used by LT (and a few test programs),
amounting to several dozen of the approximately 150 defined in the
manual.

A complete run of LT---attempting proofs of 23 theorems---takes
approximately 574,000 IPL machine cycles and finishes in a few seconds
on modern hardware. On the JOHNNIAC at RAND in 1956, it likely
required hours.

%% ============================================================
\section{What the Logic Theorist Proves}
\label{sec:results}
%% ============================================================

The Logic Theorist attempts to prove theorems in propositional logic
using the five axioms of \textit{Principia Mathematica}, Chapter~2.
The notation used by the original authors is given in
Table~\ref{tab:notation}.

\begin{table}[h]
\centering
\caption{Logical notation used by the Logic Theorist}
\label{tab:notation}
\begin{tabular}{cl}
\toprule
\textbf{Symbol} & \textbf{Meaning} \\
\midrule
\texttt{I} & Implication ($\rightarrow$) \\
\texttt{V} & Disjunction ($\lor$) \\
\texttt{*} & Conjunction ($\land$) \\
\texttt{-} & Negation ($\lnot$) \\
\texttt{=} & Biconditional ($\leftrightarrow$) \\
\texttt{.=.} & Definitional equality \\
\bottomrule
\end{tabular}
\end{table}

LT employs several proof methods:

\begin{itemize}
\item \textbf{Substitution}---substitute variables in an axiom or
  previously proved theorem to match the goal.
\item \textbf{Detachment} (modus ponens)---from $A \rightarrow B$ and
  $A$, conclude $B$.
\item \textbf{Forward chaining}---apply substitution and detachment
  working forward from known theorems.
\item \textbf{Backward chaining}---work backward from the goal,
  seeking premises that would yield it.
\item \textbf{Sublevel replacement}---replace a subexpression using a
  known equivalence.
\end{itemize}

Each proof attempt is bounded by an effort limit of 20,000 IPL
machine cycles. This is a soft cap: it prevents the system from
\textit{beginning} a new subproblem exploration once exceeded, but
cannot interrupt an in-progress search.

Table~\ref{tab:results} shows the results of the reanimated LT on the
23 theorems attempted. Of these, 16 are proved and 7 are not---results
that are historically consistent with the original LT's known
behavior within its search effort limits.

\begin{table}[h]
\centering
\caption{Theorems attempted by the reanimated Logic Theorist}
\label{tab:results}
\begin{tabular}{lll}
\toprule
\textbf{Thm.} & \textbf{Formula} & \textbf{Result} \\
\midrule
2.01 & \texttt{(PI-P)I-P} & Proved \\
2.02 & \texttt{QI(PIQ)} & Proved \\
2.04 & \texttt{(PI(QIR))I(QI(PIR))} & Proved \\
2.05 & \texttt{(QIR)I((PIQ)I(PIR))} & Proved \\
2.06 & \texttt{(PIQ)I((QIR)I(PIR))} & Proved \\
2.07 & \texttt{PI(PVP)} & Proved \\
2.08 & \texttt{PIP} & Proved \\
2.10 & \texttt{-PVP} & Proved \\
2.11 & \texttt{PV-P} & Proved \\
2.12 & \texttt{PI--P} & Proved \\
2.13 & \texttt{PV---P} & Proved \\
2.14 & \texttt{--PIP} & Not proved \\
2.15 & \texttt{(-PIQ)I(-QIP)} & Not proved \\
2.20 & \texttt{PI(PVQ)} & Proved \\
2.21 & \texttt{-PI(PIQ)} & Not proved \\
2.24 & \texttt{PI(-PVQ)} & Proved \\
3.13 & \texttt{(-(P*Q))I(-PV-Q)} & Not proved \\
3.14 & \texttt{(-PV-Q)I(-(P*Q))} & Not proved \\
3.24 & \texttt{-(P*-P)} & Not proved \\
4.13 & \texttt{P=--P} & Not proved \\
4.20 & \texttt{P=P} & Proved \\
4.24 & \texttt{P=(P*P)} & Proved \\
4.25 & \texttt{P=(PVP)} & Proved \\
\bottomrule
\end{tabular}
\end{table}

A representative proof trace is shown below. Here LT proves Theorem
2.01, \texttt{(PI-P)I-P}:

\begin{lstlisting}
2.01    (PI-P)I-P

PROOF FOUND.

    GIVEN           *1.2  (AVA)IA
    SUBSTITUTION    .0    (PI-P)IP
    SUBLEVEL REPL   2.01  (PI-P)I-P
    Q.E.D.

EFFORT      LIMIT 20000   ACTUAL 5579
SUBPROBLEMS LIMIT 50      ACTUAL 1
SUBSTITUTIONS LIMIT 50    ACTUAL 2
\end{lstlisting}

The proof begins from Axiom~*1.2, \texttt{(AVA)IA}, substitutes
\texttt{-P} for \texttt{A} to obtain \texttt{(-PV-P)I-P}, then
applies sublevel replacement (using the fact that \texttt{P} implies
\texttt{PVP}) to transform the antecedent, yielding the desired
theorem.

%% ============================================================
\section{Debugging with LLMs: A Collaboration with AI to Reanimate AI}
\label{sec:debugging}
%% ============================================================

The most distinctive aspect of this project---and perhaps its most
novel methodological contribution---was the essential role played by
large language models as collaborative debugging partners. Over the
course of approximately one year, I worked closely with two LLMs:
Google's Gemini (via the AntiGravity interface) and Anthropic's Claude
(via its terminal-based coding interface, Claude Code). The complete
histories of these interactions are preserved in the project
repository.

\subsection{Teaching a Dead Language to an LLM}

The collaboration began from a position of mutual ignorance. IPL-V is
an essentially dead language: the manual is poorly OCR'd in places and
obscure even where legible; the language itself is unlike any modern
programming language; and there is no Stack Overflow, no GitHub
corpus, no community of practitioners for the LLMs to have learned
from. I was simultaneously learning the language myself while building
the interpreter, which meant that the LLMs' primary source of
information about IPL-V was my own Lisp code---which contained bugs.

This created a fascinating epistemological bootstrapping problem. The
LLMs were inferring IPL-V semantics from my interpreter, sometimes
correctly and sometimes not. When the interpreter had bugs, the LLMs
would occasionally learn incorrect semantics and propagate those
errors in their suggestions. Untangling this required constant
vigilance about what was established knowledge versus what was
inferred from possibly-buggy code.

\subsection{The Two-Unknown Problem}

A persistent challenge was what I came to think of as the two-unknown
problem: when something went wrong, was the bug in my IPL-V
interpreter, or in the transcribed LT code? Early on, I made a
methodological decision that proved essential: we would assume the LT
code was correct. This was not entirely justified---indeed, several
transcription errors were eventually found---but it was a necessary
simplifying assumption. Without it, every anomaly would have opened
two lines of investigation, making progress nearly impossible.

This same problem was not lost on the original developers. In the
McCorduck interviews, Cliff Shaw describes the particular difficulty
of debugging heuristic programs: they do not crash, they do not
produce obviously wrong output, they simply generate vast quantities
of behavior that might or might not be correct. Sixty years later,
with the same program, I faced the identical challenge.

\subsection{The Wall---and How LLMs Broke Through It}

The debugging arc had a clear inflection point. In the early phase,
bugs manifested as crashes or obviously incorrect behavior---a
J-function not implemented, a register not saved and restored, a
stack underflow. These were frustrating but tractable. A human
programmer, even working alone, could identify and fix them.

The second phase was qualitatively different. Once the interpreter
was stable enough that LT ran without crashing, the output was
thousands of lines of execution trace---machine cycles, register
states, list operations---that \textit{might} be correct. The program
was heuristic; it was supposed to explore and backtrack. Distinguishing
a correct backtrack from one caused by a subtle bug required
understanding both the IPL-V machine semantics and LT's high-level
proof strategy simultaneously, while tracking state across thousands
of steps.

I will be frank: at this point, I had effectively given up. The
cognitive load of holding the machine state in my head across
thousands of trace lines, while simultaneously questioning whether
each line reflected correct or incorrect behavior, exceeded what I
could sustain. The LLMs, however, were not inhibited by this. They
could read thousands of lines of trace, insert diagnostic prints at
ten-line intervals, and maintain the kind of tedious but precise
state-tracking that the task demanded. It was enormously repetitive
work that simultaneously required deep contextual understanding---a
combination that, ironically, is precisely where current LLMs excel.

\subsection{Simon's J's: The Breakthrough}

A critical turning point came when the LLM had accumulated enough
understanding of IPL-V---through months of working with my interpreter
and the manual---that I was able to give it a new primary source:
``Simon's J's''~\cite{bochannek2012simonsjs}, a set of original IPL-V
punched cards from 1962, preserved at the Computer History Museum.
These cards contain Simon's own implementations of the J-functions in
IPL-V assembly.

The cards had been transcribed by museum volunteers and made available
online. When I provided these transcriptions to the LLM, it was able
to compare Simon's original J-function implementations against my Lisp
versions and identify discrepancies. This was a remarkable moment: an
AI system, having learned a dead programming language primarily from a
modern interpreter that it had helped debug, was now able to read
original 1962 source code and use it to find bugs.

\textbf{[PLACEHOLDER: A specific debugging anecdote will be inserted
here, drawn from the Claude Code session logs. This will describe a
particular J-function bug---the symptom in the LT trace, the
discrepancy between the Lisp implementation and Simon's original, and
how the LLM identified and resolved it.]}

>>>>>>>>>>>>>>>>>>>>>>>>>>>>>>>>>>>>>>>>>>>>>>>>>>>>>>>>>>>

Okay, here's the anecdote: After over a year of working with LT
running on my interpreter, it was working well enough that it was
actually proving trivial (substitution only) theorems, and wasn't
crashing. However, it also wasn't really working -- it wouldn't prove
complex theorms that it shoudl have done, and would get into loops,
not recognizing when a theorem was proved. I had a very similar
feeling to Shaw's debugging description above -- hard crashes are
easier to debug than what might be misbehavior. A hard crash is almost
always traceable to some error in the runup to the crash, but
misbehavior could be due to errors anywhere in the run, and at this
point, even for small proofs (or, rather, failed proofs) it was
running tens of thousands of interpreted machine cycles, and so I
would have to analyze tens of thousands of lines of traces and
dumps. In fact, I'd almot given up.  This is where LLMs, specifically
Gemini (AntiGravity) and Claude (the desktop coding thing) came in
very handy. They are very good at micro-programming and debugging, but
poor at keeping a large project in their heads. What was needed here
was tenacious By this time I had developed a very good sense that my
IPL-V interpreter basically worked (otherwise we'd have never got
through even a dozen instrutions!), that the LT transcription was
correct, and that therefore, almost all errors were in the
J-functions, which I had written in Lisp from their descriptions in
the Newell, et al. manual. My experience was that the description the
the J-Functions in the manual were sometimes vague in subtle ways, an
this had bit me before. I was, in particular, warry of the most
complex J-functions, the generators (e.g., J17-J19, and J100) and list
manipulation functions (e.g., J74 and others). Programming with LLMs
involves a lot of teaching of the project. In this particular case, I
also had to do a lot of teaching of IPL-V, since, unlike almost any
other languge that they deal with, there is not Stack Overflow
community for IPL-V; In fact, I may be the only one on the planet who
knows the language fluently at this point! (Probably also Ed
Feigenbaum.) What the LLMs are best at is tenacity, they're willing to
read through thousands of lines of traces and dumps, and put in
breakpoints at hundreds of places, until they localize the problem;
Work that I had given up on. Gemini found the first issue after
several hours (of LLM time!) over several days, and as predicted it
was in a J-function. This made significant progress -- now LT was
generating proofs, but wasn't stopping when a proof was found. After
literally another several hours (of LLM time!) Claude (desktop)
isolated the second problem in J74 -- list copying -- but we didn't
know how to fix it because both Claude and I agreed that my code
seemed to match the description, but there was a clear problem. This
is where it actually gets interesting. We had Simon's Js from the
Computer History Museum, and I recalled the J74 was among them. By now
I'd taught Claude enough IPL-V that it was able to read the IPL-V code
in Simon's Js, and compare it with my Lisp implemention of J74, and
guess what ... they differed in exactly the hypothesized way! A one
line patch to my J74 implementation made LT run completely correctly!

<<<<<<<<<<<<<<<<<<<<<<<<<<<<<<<<<<<<<<<<<<<<<<<<<<<<<<<<<<<

\subsection{A Closer Collaboration Than Any Human}

One observation deserves explicit mention: over the course of this
project, my collaboration with the LLMs was closer than with any human
engineer. This is not a slight to the human collaborators, whose
contributions---particularly in transcription and encouragement---were
essential. But the day-to-day, line-by-line work of understanding
IPL-V semantics, reading traces, forming and testing hypotheses about
bugs, and iterating on fixes was done in dialogue with Claude and
Gemini. The LLMs served simultaneously as research assistants,
sounding boards, and---once they had absorbed enough context---as
something approaching genuine collaborators who understood the domain.

Both LLMs made substantive contributions. Gemini found one major bug;
Claude found several others. In the end, I settled on working
primarily with Claude, not because of a significant capability
difference, but because it was easier to maintain a single long
session with accumulated context than to switch between systems.

%% ============================================================
\section{Reflections: IPL, Lisp, and the Arc of AI}
\label{sec:reflections}
%% ============================================================

I have been programming in Lisp for fifty years. For most of that
time, I attributed the genius of Lisp entirely to John McCarthy. One
of the most valuable outcomes of this project has been a revision of
that understanding.

Building an IPL-V interpreter revealed that Newell, Simon, and Shaw
had, by the late 1950s, already invented essentially every conceptual
element that makes Lisp powerful: list manipulation, symbol processing,
the idea that programs are themselves lists, dynamic memory allocation,
recursion, higher-order functions, generators. What they lacked was
the all-important homoiconic syntax---the unification of code and data
in a single parenthesized notation that makes Lisp programs
manipulable by Lisp programs.

This is not to diminish McCarthy's contribution, which was genuinely
brilliant. But it reframes the standard narrative. McCarthy's genius
was not in inventing list processing \textit{ex nihilo}, but in
finding the syntactic form that made these ideas---ideas that Newell,
Simon, and Shaw had already developed---elegant and accessible. IPL-V
and Lisp coexisted for approximately five years, but Lisp was so much
simpler and more expressive that it rapidly supplanted its ancestor.

There is a personal dimension to this realization. I studied under
both Newell and Simon at Carnegie Mellon University in the 1980s. IPL
and LT were long since history by that time, and I never discussed
either with them---we barely learned about these pioneering efforts in
our classes. But something from working with these giants stuck with
me. At CMU I was modeling scientific and engineering reasoning, and
continued to do so for most of my career. Now, at the end of that
career, I have been led back to try to understand how these brilliant
scientist-engineers were inventing AI and cognitive science.

%% ============================================================
\section{Related Work}
\label{sec:related}
%% ============================================================

The reanimation of historical software systems has gained momentum in
recent years. Most directly related is the ELIZA reanimation by Lane
et al.~\cite{lane2025eliza}, in which a team including the present
author restored Weizenbaum's original ELIZA to operation on a
reconstructed CTSS running on an emulated IBM~7094. That project
reconstructed the complete original hardware and operating system
stack; the present work instead emulates the abstract machine
specification in a modern language.

A 1968 paper describes a reimplementation of the Logic Theorist in
LISP~1.5~\cite{geiger1968lt}. That effort, however, rewrote LT's
logic in Lisp rather than faithfully emulating IPL-V and running the
original code. The distinction is significant: a Lisp reimplementation
demonstrates that LT's algorithms can be expressed in Lisp (which is
unsurprising), while the present work demonstrates that the original
IPL-V code, as written by Newell, Simon, and Shaw's team, still
functions correctly when provided with a faithful implementation of
the machine it was designed for.

To the author's knowledge, no prior attempt has been made to build an
IPL-V interpreter and execute original IPL-V programs.

%% ============================================================
\section{Next Steps}
\label{sec:next}
%% ============================================================

The IPL-V interpreter developed for this project is not specific to
the Logic Theorist. It implements the general IPL-V abstract machine
and can, in principle, run any IPL-V program. The most significant
targets for future reanimation are the General Problem Solver
(GPS)~\cite{newell1959gps} and EPAM (Elementary Perceiver and
Memorizer)~\cite{feigenbaum1961epam}, both originally written in
IPL-V by members of the same research group. Archival work to locate
complete source code for GPS is planned for spring 2026. Now that a
proven IPL-V interpreter exists, running these programs should require
substantially less effort than the initial construction of the
interpreter itself.

The interpreter, the transcribed LT code, and the complete LLM
interaction histories are available as open source at
\url{https://github.com/jeffshrager/IPL-V}.

%% ============================================================
\section{Acknowledgements}
%% ============================================================

The following people contributed to this project in ways ranging from
code transcription to weekly encouragement: Anthony Hay, Arthur
Schwarz, Leigh Klotz, David M.\ Berry, Rupert Lane, Amy Majczyk,
Ethan Ableman, and Paul McJones. I am particularly grateful to Hay,
Schwarz, Klotz, and Berry, who received weekly progress updates
throughout the project and responded with encouragement that was
always helpful and sometimes technically substantive. Rupert Lane's
gridlock software caught several transcription errors that had
survived multiple rounds of human review.

I also thank the Computer History Museum for preserving and making
available ``Simon's J's,'' the original 1962 IPL-V J-function card
deck, which proved essential during debugging.

Finally, I want to acknowledge Allen Newell and Herbert Simon, under
whom I studied at CMU in the 1980s. We never discussed IPL or the
Logic Theorist, but the intellectual values they instilled---rigor
about mechanism, respect for the complexity of cognition, and the
conviction that thinking could be understood computationally---led me,
decades later, back to their earliest work.

\begin{thebibliography}{99}


>>>>>>>>>>>>>>>>>>>>>>>>>>>>>>>>>>>>>>>>>>>>>>>>>>>>>>>>>>>

@article{richman1995epam,
  author  = {Richman, Howard B. and Staszewski, James J. and Simon, Herbert A.},
  title   = {Simulation of Expert Memory Using {EPAM} {IV}},
  journal = {Psychological Review},
  volume  = {102},
  number  = {2},
  pages   = {305--330},
  year    = {1995},
  doi     = {10.1037/0033-295X.102.2.305}
}

@book{whitehead1910principia,
  author    = {Alfred North Whitehead and Bertrand Russell},
  title     = {Principia Mathematica},
  publisher = {Cambridge University Press},
  address   = {Cambridge},
  year      = {1910},
  note      = {3 volumes published 1910--1913}
}

@book{newell1961ipl,
  editor    = {Allen Newell},
  title     = {Information Processing Language-V Manual},
  publisher = {Prentice-Hall},
  address   = {Englewood Cliffs, NJ},
  year      = {1961}
}

@article{mccarthy1960lisp,
  author  = {John McCarthy},
  title   = {Recursive Functions of Symbolic Expressions and Their Computation by Machine, Part I},
  journal = {Communications of the ACM},
  volume  = {3},
  number  = {4},
  pages   = {184--195},
  month   = apr,
  year    = {1960}
}

@techreport{stefferud1963lt,
  author      = {Erik Stefferud},
  title       = {The Logic Theory Machine: A Model Heuristic Program},
  institution = {RAND Corporation},
  address     = {Santa Monica, CA},
  number      = {RM-3731-CC},
  year        = {1963}
}

@book{newell1964iplv,
  editor    = {Allen Newell},
  title     = {Information Processing Language-V Manual},
  edition   = {2},
  publisher = {Prentice-Hall},
  address   = {Englewood Cliffs, NJ},
  year      = {1964}
}

@article{lane2025eliza,
  author  = {Rupert Lane and Alexander Hay and Andreas Schwarz and David M. Berry and Jeff Shrager},
  title   = {ELIZA Reanimated: Restoring the Mother of All Chatbots to One of the World's First Time-Sharing Systems},
  journal = {IEEE Annals of the History of Computing},
  volume  = {47},
  number  = {1},
  pages   = {8--23},
  year    = {2025}
}

@techreport{gruenberger1968johnniac,
  author      = {Fred Gruenberger},
  title       = {The JOHNNIAC},
  institution = {RAND Corporation},
  address     = {Santa Monica, CA},
  number      = {RM-5654-PR},
  year        = {1968}
}

@book{mccorduck1979machines,
  author    = {Pamela McCorduck},
  title     = {Machines Who Think},
  publisher = {W. H. Freeman},
  address   = {San Francisco},
  year      = {1979}
}

@misc{bochannek2012simonsjs,
  author       = {Andrew Bochannek},
  title        = {Simon's J's},
  howpublished = {Computer History Museum Blog},
  year         = {2012},
  month        = sep,
  day          = {26},
  url          = {https://computerhistory.org/blog/simons-js/}
}

@misc{lane2024gridlock,
  author = {Rupert Lane},
  title  = {Gridlock: Accurate Code Extraction from PDFs},
  year   = {2024},
  url    = {https://github.com/rupertl/gridlock},
  note   = {GitHub repository}
}

@article{millstein1968logic,
  author  = {J. S. Millstein},
  title   = {The Logic Theorist in LISP},
  journal = {The Computer Journal},
  volume  = {11},
  number  = {1},
  pages   = {39--44},
  month   = feb,
  year    = {1968}
}

@inproceedings{newell1959gps,
  author    = {Allen Newell and J. C. Shaw and Herbert A. Simon},
  title     = {Report on a General Problem-Solving Program},
  booktitle = {Proceedings of the International Conference on Information Processing},
  address   = {Paris, France},
  pages     = {256--264},
  year      = {1959}
}

@inproceedings{feigenbaum1961epam,
  author    = {Edward A. Feigenbaum},
  title     = {The Simulation of Verbal Learning Behavior},
  booktitle = {Proceedings of the Western Joint Computer Conference},
  pages     = {121--132},
  year      = {1961}
}end{thebibliography}

<<<<<<<<<<<<<<<<<<<<<<<<<<<<<<<<<<<<<<<<<<<<<<<<<<<<<<<<<<<

\end{document}


Some notes for thing to add:

-- Although the Logic Theorist and ELIZA addressed very different
problems—formal theorem proving and natural-language conversation—both
programs emerged from the same symbolic programming lineage
originating in IPL and later realized in Lisp and related
list-processing systems. Each can be understood as a rule-based
control system operating over symbolic list structures.

-- The famous “shorter proof than Principia”

One of Logic Theorist’s results is historically important because it improved on Whitehead & Russell.
The system produced a shorter proof of Theorem 2.85 in Principia Mathematica.
Whitehead and Russell’s proof:

  long derivation using intermediate lemmas

Logic Theorist:

  found a more direct proof path

Simon later wrote:

  “We discovered that the machine had found a proof of one theorem that was more elegant than the one in Principia.”

He joked that if they had submitted it:

  “Russell might have been pleased.”

This anecdote appears in several later retrospectives but is rarely
included in reconstruction papers, even though it illustrates the
heuristic search idea beautifully.

You could mention this in the background section when explaining LT’s significance.

-- The program was originally hand-simulated

Before the program ran on JOHNNIAC, Newell and Simon simulated parts of it by hand.

They literally:

  wrote expressions on paper
  applied IPL-like list operations manually
  tracked stacks and subgoals

Simon described this as: “running the program on the human computer.”

This is historically fascinating because it mirrors Turing’s original
machine simulations and later hand-simulations of early neural
networks.

You could mention this when discussing the difficulty of implementing LT.

-- The heuristics were inspired by human problem-solving studies

Logic Theorist was not just a theorem prover. The heuristics came from protocol analysis of human reasoning.

Two major ideas:

1.  Means–ends analysis: Choose operations that reduce the difference between the current state and a goal state

This later becomes the core of GPS (General Problem Solver).

2. Subgoal decomposition: Break the proof goal into smaller lemmas.

These ideas become central to cognitive architectures:

  GPS
  SOAR
  ACT

EPAM-style memory models:

So LT sits right at the beginning of AI + cognitive psychology fusion.

Given your EPAM citations, this connection is especially relevant.

A nice one-sentence line you could add

Something like:

``Beyond reproducing Principia proofs, the Logic Theorist introduced
heuristic search strategies—means–ends analysis and subgoal
decomposition—that later became foundational ideas in both artificial
intelligence and cognitive architectures such as GPS and EPAM.''

That line would connect your bibliography very nicely.

Bonus obscure detail (reviewers love this):

  The first LT run produced 52 theorems overnight on the JOHNNIAC.

Simon reportedly walked into Newell’s office and said:

  “Over Christmas, Al, we invented thinking.”

It's a famous anecdote, but still delightful.
